\documentclass[11pt,a4paper]{article}
\usepackage[utf8]{inputenc}
\usepackage[T1]{fontenc}
\usepackage[french]{babel}
\usepackage{lmodern}
\usepackage[margin=1.8cm]{geometry}
\usepackage{amsmath,amssymb}
\usepackage{array}
\usepackage{enumitem}
\usepackage{xcolor}
\usepackage[most]{tcolorbox}
\usepackage{tikz}
\usetikzlibrary{arrows.meta}

\setlength{\parindent}{0pt}
\setlength{\parskip}{4pt}
\renewcommand{\arraystretch}{1.35}

\newcommand{\R}{\mathbb{R}}

\newtcolorbox{objbox}{colback=blue!4,colframe=blue!45!black,boxrule=0.6pt,arc=2pt,
  left=6pt,right=6pt,top=5pt,bottom=5pt}
\newtcolorbox{retenir}{colback=orange!6,colframe=orange!75!black,boxrule=0.6pt,arc=2pt,
  left=6pt,right=6pt,top=5pt,bottom=5pt,title={\bfseries À retenir},fonttitle=\small}
\newtcolorbox{corrbox}{colback=green!3,colframe=green!45!black,boxrule=0.6pt,arc=2pt,
  left=6pt,right=6pt,top=5pt,bottom=5pt}

\newcommand{\partie}[2]{%
  \vspace{6pt}
  \noindent\colorbox{blue!12}{\parbox{\dimexpr\linewidth-2\fboxsep}{%
    \bfseries\color{blue!50!black} #1 \hfill \normalfont\small\itshape (durée indicative : #2)}}%
  \vspace{4pt}}

\begin{document}

% ===================== EN-TÊTE =====================
\begin{center}
  {\Large\bfseries Activité — Le second degré, partie 1}\\[2pt]
  {\large De la fonction carrée à la forme canonique}\\[6pt]
  {\small Première STI2D · Mathématiques · Durée : 50 minutes}\\[6pt]
  \begin{tabular}{@{}p{6cm}p{6cm}@{}}
    Nom : \dotfill & Prénom : \dotfill \\[8pt]
    Classe : \dotfill & Date : \dotfill \\
  \end{tabular}
\end{center}

\begin{objbox}
\textbf{Objectifs.} 
\begin{itemize}[leftmargin=1.2em,itemsep=1pt,topsep=2pt]
  \item Rappeler les propriétés de la fonction carrée : parité, variations, signe, limites.
  \item Reconnaître l'effet d'une translation horizontale, verticale puis oblique sur une parabole.
  \item Reconnaître l'effet d'une affinité (coefficient $k$) : ouverture, retournement.
  \item Lire sur la forme canonique $k(x-\alpha)^2+\beta$ : sommet, axe, extremum, variations, signe et limites.
  \item Calculer des images et contrôler la cohérence d'un résultat.
\end{itemize}
\textbf{Matériel.} Le document interactif \texttt{DEGRE2.html} est projeté au tableau : il sert à vérifier chaque réponse (curseurs, tableaux de variations et de signes, calculs d'images).
\end{objbox}

% ===================== PARTIE A =====================
\partie{Partie A — Rappels : la fonction carrée}{5 min}

Soit $f$ la fonction définie sur $\R$ par $f(x)=x^2$, et $\mathcal{P}$ sa courbe dans un repère.

\begin{enumerate}[leftmargin=1.4em,itemsep=3pt]
  \item Compléter le tableau de valeurs :
  \begin{center}
  \begin{tabular}{|c|c|c|c|c|c|c|c|}
  \hline
  $x$ & $-3$ & $-2$ & $-1$ & $0$ & $1$ & $2$ & $3$\\ \hline
  $f(x)=x^2$ & & & & & & & \\ \hline
  \end{tabular}
  \end{center}

  \item Justifier que, pour tout réel $x$, $f(-x)=f(x)$. \\
  Que peut-on en déduire pour la courbe $\mathcal{P}$ ? Comment s'appelle cette propriété ?

  \item Compléter le tableau de variations de $f$ :
  \begin{center}
  $\begin{array}{c|ccccc}
  \hline
  x & -\infty & & 0 & & +\infty\\ \hline
  f(x) & & & & & \\ \hline
  \end{array}$
  \end{center}

  \item Compléter le tableau de signes de $f$ :
  \begin{center}
  $\begin{array}{c|ccccc}
  \hline
  x & -\infty & & 0 & & +\infty\\ \hline
  f(x) & & & & & \\ \hline
  \end{array}$
  \end{center}

  \item Écrire les deux limites de $f$ en $-\infty$ et en $+\infty$.

  \item Vrai ou faux ? Justifier chaque réponse.
  \begin{enumerate}[label=\alph*),leftmargin=1.6em,itemsep=1pt]
    \item « Pour tout réel $x$, on a $x^2>0$. »
    \item « L'équation $x^2=9$ admet exactement deux solutions. »
    \item « $f(-5)=f(5)$. »
  \end{enumerate}
\end{enumerate}

% ===================== PARTIE B =====================
\partie{Partie B — Translation horizontale}{8 min}

Soit $g$ la fonction définie sur $\R$ par $g(x)=(x-2)^2$.

\begin{enumerate}[leftmargin=1.4em,itemsep=3pt]
  \item Calculer $g(0)$, $g(1)$, $g(2)$ et $g(3)$.

  \item En remarquant que $g(x)=f(x-2)$, décrire la transformation qui permet d'obtenir la courbe de $g$
  à partir de celle de $f$. Donner le vecteur de translation.

  \item Donner les coordonnées du sommet $S$ de la parabole de $g$ et l'équation de son axe de symétrie.

  \item Dresser le tableau de variations de $g$.

  \item Dresser le tableau de signes de $g$.

  \item Écrire les limites de $g$ en $-\infty$ et en $+\infty$.

  \item Sans faire de calcul, expliquer pourquoi $g(-1)=g(5)$.
\end{enumerate}

\begin{center}
\begin{tikzpicture}[scale=0.55,>=Stealth]
  \draw[gray!22,very thin,step=1] (-3,0) grid (6,9);
  \draw[->,thick] (-3,0) -- (6.6,0) node[below right] {$x$};
  \draw[->,thick] (0,-0.3) -- (0,9.6) node[left] {$y$};
  \foreach \x in {-3,-2,-1,1,2,3,4,5,6} \draw (\x,0.15) -- (\x,-0.15) node[below,font=\scriptsize]{\x};
  \foreach \y in {1,2,...,9} \draw (0.15,\y) -- (-0.15,\y) node[left,font=\scriptsize]{\y};
  \node[below left,font=\scriptsize] at (0,0) {$O$};
  \draw[gray,dashed,domain=-3:3,samples=90] plot (\x,{\x*\x});
  \draw[blue,thick,domain=-1:5,samples=120] plot (\x,{(\x-2)*(\x-2)});
  \fill[red] (2,0) circle (2.6pt);
  \node[red,above right,font=\scriptsize] at (2,0) {$S(2;0)$};
\end{tikzpicture}\\[-2pt]
{\small Figure 1 — Courbe de $g(x)=(x-2)^2$ (en bleu) et courbe de référence $y=x^2$ (en pointillés).}
\end{center}

% ===================== PARTIE C =====================
\partie{Partie C — Translation verticale}{7 min}

Soit $h$ la fonction définie sur $\R$ par $h(x)=x^2-3$.

\begin{enumerate}[leftmargin=1.4em,itemsep=3pt]
  \item Décrire la transformation qui permet d'obtenir la courbe de $h$ à partir de celle de $f$.
  Donner le vecteur de translation.

  \item Donner les coordonnées du sommet et l'équation de l'axe de symétrie.

  \item Résoudre $h(x)=0$. Donner les valeurs exactes puis des valeurs arrondies au centième.

  \item Dresser le tableau de signes de $h$.

  \item Dresser le tableau de variations de $h$ et préciser le minimum.

  \item Écrire les limites de $h$ en $-\infty$ et en $+\infty$.

  \item \textbf{Synthèse.} On considère $x\mapsto x^2+c$ où $c$ est un réel quelconque.
  Discuter, suivant le signe de $c$, le nombre de solutions de l'équation $x^2+c=0$
  et le signe de $x^2+c$ sur $\R$.
\end{enumerate}

% ===================== PARTIE D =====================
\partie{Partie D — Translation oblique}{8 min}

Soit $j$ la fonction définie sur $\R$ par $j(x)=(x+1)^2-4$.

\begin{enumerate}[leftmargin=1.4em,itemsep=3pt]
  \item Donner les coordonnées du sommet, l'équation de l'axe de symétrie et le vecteur de la translation
  qui transforme la parabole de référence en celle de $j$.

  \item Résoudre $j(x)=0$ \emph{(indication : commencer par écrire $(x+1)^2=4$)}.

  \item Dresser le tableau de signes de $j$.

  \item Dresser le tableau de variations de $j$.

  \item Écrire les limites de $j$ en $-\infty$ et en $+\infty$.

  \item Calculer $j(0)$, $j(2)$ et $j(-4)$. Que remarque-t-on pour $j(2)$ et $j(-4)$ ?
  Faire le lien avec l'axe de symétrie.

  \item Lire sur la figure 2 les solutions de $j(x)=0$, puis vérifier avec la question 2.
\end{enumerate}

\begin{center}
\begin{tikzpicture}[scale=0.6,>=Stealth]
  \draw[gray!22,very thin,step=1] (-5,-5) grid (4,6);
  \draw[->,thick] (-5,0) -- (4.6,0) node[below right] {$x$};
  \draw[->,thick] (0,-5) -- (0,6.6) node[left] {$y$};
  \foreach \x in {-5,-4,-3,-2,-1,1,2,3,4} \draw (\x,0.15) -- (\x,-0.15) node[below,font=\scriptsize]{\x};
  \foreach \y in {-5,-4,-3,-2,-1,1,2,3,4,5,6} \draw (0.15,\y) -- (-0.15,\y) node[left,font=\scriptsize]{\y};
  \node[below left,font=\scriptsize] at (0,0) {$O$};
  \draw[gray,dashed,domain=-2.4:2.4,samples=90] plot (\x,{\x*\x});
  \draw[blue,thick,domain=-4.1:2.1,samples=130] plot (\x,{(\x+1)*(\x+1)-4});
  \fill[red] (-1,-4) circle (2.6pt);
  \node[red,below right,font=\scriptsize] at (-1,-4) {$S(-1;-4)$};
\end{tikzpicture}\\[-2pt]
{\small Figure 2 — Courbe de $j(x)=(x+1)^2-4$ (en bleu) et courbe de référence $y=x^2$ (en pointillés).}
\end{center}

% ===================== PARTIE E =====================
\partie{Partie E — Affinité (dilatation verticale)}{7 min}

On considère les fonctions $k_1(x)=2x^2$ et $k_2(x)=-0{,}5\,x^2$.

\begin{enumerate}[leftmargin=1.4em,itemsep=3pt]
  \item Compléter le tableau de valeurs :
  \begin{center}
  \begin{tabular}{|c|c|c|c|c|c|}
  \hline
  $x$ & $-2$ & $-1$ & $0$ & $1$ & $2$\\ \hline
  $k_1(x)=2x^2$ & & & & & \\ \hline
  $k_2(x)=-0{,}5\,x^2$ & & & & & \\ \hline
  \end{tabular}
  \end{center}

  \item Dresser les tableaux de variations de $k_1$ et de $k_2$. Préciser pour chacune la nature
  de l'extremum (minimum ou maximum) et sa valeur.

  \item Donner le tableau de signes de $k_1$, puis celui de $k_2$.

  \item Écrire les limites de $k_1$ et de $k_2$ en $-\infty$ et en $+\infty$.

  \item Que change le \emph{signe} de $k$ sur l'allure de la parabole ? Que change sa \emph{valeur absolue} ?
\end{enumerate}

% ===================== PARTIE F =====================
\partie{Partie F — Synthèse : la forme canonique}{10 min}

Soit $f$ la fonction définie sur $\R$ par
\[
  f(x)=-2(x-3)^2+8 .
\]

\begin{enumerate}[leftmargin=1.4em,itemsep=3pt]
  \item Identifier les réels $k$, $\alpha$ et $\beta$ tels que $f(x)=k(x-\alpha)^2+\beta$.

  \item En déduire les coordonnées du sommet $S$ de la parabole, l'équation de son axe de symétrie,
  ainsi que la nature et la valeur de l'extremum.

  \item Dresser le tableau de variations de $f$.

  \item Résoudre $f(x)=0$, puis dresser le tableau de signes de $f$.

  \item Écrire les limites de $f$ en $-\infty$ et en $+\infty$.

  \item Calculer $f(0)$, $f(3)$, $f(4)$ et $f(2)$. Vérifier que $f(4)=f(2)$ et expliquer ce résultat
  par la symétrie de la parabole.

  \item Les points $A(0\,;-10)$ et $B(1\,;0)$ appartiennent-ils à la courbe de $f$ ? Justifier par un calcul.

  \item \textbf{Bilan.} Rédiger en trois lignes maximum la méthode de lecture de la forme canonique :
  que lit-on sur $k$ ? sur $\alpha$ ? sur $\beta$ ?
\end{enumerate}

\begin{retenir}
Avec $f(x)=k(x-\alpha)^2+\beta$ et $k\neq 0$ :
sommet $S(\alpha;\beta)$ · axe de symétrie $x=\alpha$ · extremum $\beta$ atteint en $\alpha$ ·
si $k>0$, c'est un minimum (parabole « vers le haut ») ; si $k<0$, un maximum (parabole « retournée ») ·
le signe de $f$ se lit à partir de la position de $\beta$ par rapport à $0$ et du signe de $k$ ·
les limites en $\pm\infty$ ont le signe de $k$.
\end{retenir}

\vfill
{\small\itshape Prolongement (partie 2) : forme développée $ax^2+bx+c$, factorisation, résolution des équations
du second degré et signe du trinôme.}

% ===================== CORRIGÉ =====================
\newpage
\begin{center}
  {\Large\bfseries CORRIGÉ — Activité « Le second degré, partie 1 »}
\end{center}

\partie{Corrigé A — Rappels}{}

\begin{enumerate}[leftmargin=1.4em,itemsep=3pt]
  \item $f(-3)=9$, $f(-2)=4$, $f(-1)=1$, $f(0)=0$, $f(1)=1$, $f(2)=4$, $f(3)=9$.

  \item $f(-x)=(-x)^2=x^2=f(x)$. La courbe est donc symétrique par rapport à l'axe des ordonnées :
  la fonction $f$ est \textbf{paire}.

  \item Tableau de variations :
  \begin{center}
  $\begin{array}{c|ccccc}
  \hline
  x & -\infty & & 0 & & +\infty\\ \hline
  f(x) & +\infty & \searrow & 0 & \nearrow & +\infty\\ \hline
  \end{array}$
  \end{center}
  $f$ est décroissante sur $]-\infty;0]$ et croissante sur $[0;+\infty[$ ; minimum $0$ atteint en $0$.

  \item Tableau de signes :
  \begin{center}
  $\begin{array}{c|ccccc}
  \hline
  x & -\infty & & 0 & & +\infty\\ \hline
  f(x) & & + & 0 & + & \\ \hline
  \end{array}$
  \end{center}
  $x^2\geq 0$ pour tout réel $x$, et $x^2=0$ seulement pour $x=0$.

  \item $\displaystyle\lim_{x\to-\infty}x^2=+\infty$ et $\displaystyle\lim_{x\to+\infty}x^2=+\infty$.

  \item
  \begin{enumerate}[label=\alph*),leftmargin=1.6em,itemsep=1pt]
    \item \textbf{Faux} : $x^2\geq 0$ et $x^2=0$ lorsque $x=0$. L'inégalité est large, pas stricte.
    \item \textbf{Vrai} : $x^2=9\iff x=3$ ou $x=-3$ (deux solutions opposées).
    \item \textbf{Vrai} : $(-5)^2=25=5^2$ ; c'est la parité de la fonction carrée.
  \end{enumerate}
\end{enumerate}

\partie{Corrigé B — Translation horizontale}{}

\begin{enumerate}[leftmargin=1.4em,itemsep=3pt]
  \item $g(0)=(-2)^2=4$ ; $g(1)=(-1)^2=1$ ; $g(2)=0$ ; $g(3)=1$.

  \item $g(x)=f(x-2)$ : la courbe de $g$ est l'image de celle de $f$ par la translation de vecteur
  $\vec{u}(2\,;0)$ (glissement de $2$ unités vers la droite).

  \item Sommet $S(2\,;0)$ ; axe de symétrie : la droite d'équation $x=2$.

  \item Tableau de variations :
  \begin{center}
  $\begin{array}{c|ccccc}
  \hline
  x & -\infty & & 2 & & +\infty\\ \hline
  g(x) & +\infty & \searrow & 0 & \nearrow & +\infty\\ \hline
  \end{array}$
  \end{center}
  Minimum $0$ atteint en $x=2$.

  \item Tableau de signes :
  \begin{center}
  $\begin{array}{c|ccccc}
  \hline
  x & -\infty & & 2 & & +\infty\\ \hline
  g(x) & & + & 0 & + & \\ \hline
  \end{array}$
  \end{center}
  $g$ est positive sur $\R$ et s'annule uniquement en $2$.

  \item $\displaystyle\lim_{x\to-\infty}g(x)=+\infty$ et $\displaystyle\lim_{x\to+\infty}g(x)=+\infty$.

  \item $-1$ et $5$ sont symétriques par rapport à $2$ (leur moyenne vaut $2$). La parabole ayant pour axe
  $x=2$, deux abscisses symétriques ont la même image : $g(-1)=g(5)=9$.
\end{enumerate}

\partie{Corrigé C — Translation verticale}{}

\begin{enumerate}[leftmargin=1.4em,itemsep=3pt]
  \item La courbe de $h$ est l'image de celle de $f$ par la translation de vecteur $\vec{u}(0\,;-3)$
  (glissement de $3$ unités vers le bas).

  \item Sommet $S(0\,;-3)$ ; axe de symétrie : l'axe des ordonnées ($x=0$).

  \item $h(x)=0\iff x^2=3\iff x=-\sqrt{3}$ ou $x=\sqrt{3}$.
  Valeurs exactes : $-\sqrt{3}$ et $\sqrt{3}$ ; valeurs arrondies : $-1{,}73$ et $1{,}73$.

  \item Tableau de signes :
  \begin{center}
  $\begin{array}{c|ccccccc}
  \hline
  x & -\infty & & -\sqrt{3} & & \sqrt{3} & & +\infty\\ \hline
  h(x) & & + & 0 & - & 0 & + & \\ \hline
  \end{array}$
  \end{center}
  $h$ est positive à l'extérieur des racines, négative entre elles.

  \item Tableau de variations :
  \begin{center}
  $\begin{array}{c|ccccc}
  \hline
  x & -\infty & & 0 & & +\infty\\ \hline
  h(x) & +\infty & \searrow & -3 & \nearrow & +\infty\\ \hline
  \end{array}$
  \end{center}
  Minimum $-3$ atteint en $x=0$.

  \item $\displaystyle\lim_{x\to-\infty}h(x)=+\infty$ et $\displaystyle\lim_{x\to+\infty}h(x)=+\infty$.

  \item \textbf{Synthèse.}
  \begin{itemize}[leftmargin=1.2em,itemsep=1pt]
    \item Si $c>0$ : $x^2+c\geq c>0$, donc \textbf{aucune solution} et $x^2+c>0$ sur $\R$.
    \item Si $c=0$ : $x^2=0$ a pour seule solution $0$, et $x^2\geq 0$ sur $\R$.
    \item Si $c<0$ : deux solutions $x=\pm\sqrt{-c}$ ; $x^2+c<0$ entre ces racines, $x^2+c>0$ à l'extérieur.
  \end{itemize}
\end{enumerate}

\partie{Corrigé D — Translation oblique}{}

\begin{enumerate}[leftmargin=1.4em,itemsep=3pt]
  \item Sommet $S(-1\,;-4)$ ; axe de symétrie $x=-1$ ; translation de vecteur $\vec{u}(-1\,;-4)$
  (le sommet passe de $O(0;0)$ à $S(-1;-4)$).

  \item $j(x)=0\iff (x+1)^2=4\iff x+1=2$ ou $x+1=-2\iff x=1$ ou $x=-3$.
  Les solutions sont donc $-3$ et $1$.

  \item Tableau de signes :
  \begin{center}
  $\begin{array}{c|ccccccc}
  \hline
  x & -\infty & & -3 & & 1 & & +\infty\\ \hline
  j(x) & & + & 0 & - & 0 & + & \\ \hline
  \end{array}$
  \end{center}

  \item Tableau de variations :
  \begin{center}
  $\begin{array}{c|ccccc}
  \hline
  x & -\infty & & -1 & & +\infty\\ \hline
  j(x) & +\infty & \searrow & -4 & \nearrow & +\infty\\ \hline
  \end{array}$
  \end{center}
  Minimum $-4$ atteint en $x=-1$.

  \item $\displaystyle\lim_{x\to-\infty}j(x)=+\infty$ et $\displaystyle\lim_{x\to+\infty}j(x)=+\infty$.

  \item $j(0)=1-4=-3$ ; $j(2)=9-4=5$ ; $j(-4)=9-4=5$.
  On a $j(2)=j(-4)$ : les abscisses $2$ et $-4$ sont symétriques par rapport à $-1$
  (moyenne $\frac{2+(-4)}{2}=-1$), axe de symétrie de la parabole.

  \item Sur la figure 2, la parabole coupe l'axe des abscisses en $x=-3$ et $x=1$, ce qui confirme la question 2.
\end{enumerate}

\partie{Corrigé E — Affinité}{}

\begin{enumerate}[leftmargin=1.4em,itemsep=3pt]
  \item
  \begin{center}
  \begin{tabular}{|c|c|c|c|c|c|}
  \hline
  $x$ & $-2$ & $-1$ & $0$ & $1$ & $2$\\ \hline
  $k_1(x)=2x^2$ & $8$ & $2$ & $0$ & $2$ & $8$\\ \hline
  $k_2(x)=-0{,}5\,x^2$ & $-2$ & $-0{,}5$ & $0$ & $-0{,}5$ & $-2$\\ \hline
  \end{tabular}
  \end{center}

  \item
  \begin{center}
  $k_1:\ \begin{array}{c|ccccc}
  \hline
  x & -\infty & & 0 & & +\infty\\ \hline
  k_1(x) & +\infty & \searrow & 0 & \nearrow & +\infty\\ \hline
  \end{array}$
  \hspace{1.2cm}
  $k_2:\ \begin{array}{c|ccccc}
  \hline
  x & -\infty & & 0 & & +\infty\\ \hline
  k_2(x) & -\infty & \nearrow & 0 & \searrow & -\infty\\ \hline
  \end{array}$
  \end{center}
  $k_1$ admet un \textbf{minimum} égal à $0$ en $0$ ; $k_2$ admet un \textbf{maximum} égal à $0$ en $0$.

  \item $k_1(x)=2x^2\geq 0$, nul en $0$ seulement ; $k_2(x)=-0{,}5\,x^2\leq 0$, nul en $0$ seulement.

  \item $\displaystyle\lim_{x\to\pm\infty}k_1(x)=+\infty$ et $\displaystyle\lim_{x\to\pm\infty}k_2(x)=-\infty$.

  \item Le \textbf{signe} de $k$ donne l'orientation : parabole tournée vers le haut (minimum) si $k>0$,
  retournée vers le bas (maximum) si $k<0$. La \textbf{valeur absolue} de $k$ règle l'ouverture :
  la parabole se resserre si $|k|>1$, s'élargit si $0<|k|<1$. Le sommet reste $O(0;0)$.
\end{enumerate}

\partie{Corrigé F — Forme canonique}{}

\begin{enumerate}[leftmargin=1.4em,itemsep=3pt]
  \item $f(x)=-2(x-3)^2+8$, donc $k=-2$, $\alpha=3$ et $\beta=8$.

  \item Sommet $S(3\,;8)$ ; axe de symétrie $x=3$ ; comme $k<0$, le sommet est un \textbf{maximum}
  de valeur $8$, atteint en $x=3$.

  \item Tableau de variations :
  \begin{center}
  $\begin{array}{c|ccccc}
  \hline
  x & -\infty & & 3 & & +\infty\\ \hline
  f(x) & -\infty & \nearrow & 8 & \searrow & -\infty\\ \hline
  \end{array}$
  \end{center}
  $f$ est croissante sur $]-\infty;3]$ et décroissante sur $[3;+\infty[$.

  \item $f(x)=0\iff -2(x-3)^2+8=0\iff (x-3)^2=4\iff x-3=2$ ou $x-3=-2$,
  c'est-à-dire $x=5$ ou $x=1$.
  \begin{center}
  $\begin{array}{c|ccccccc}
  \hline
  x & -\infty & & 1 & & 5 & & +\infty\\ \hline
  f(x) & & - & 0 & + & 0 & - & \\ \hline
  \end{array}$
  \end{center}
  $f$ est positive entre les racines et négative à l'extérieur (car $k<0$).

  \item $\displaystyle\lim_{x\to-\infty}f(x)=-\infty$ et $\displaystyle\lim_{x\to+\infty}f(x)=-\infty$
  (limites du signe de $k$).

  \item $f(0)=-2\times 9+8=-10$ ; $f(3)=8$ ; $f(4)=-2\times 1+8=6$ ; $f(2)=-2\times 1+8=6$.
  On a bien $f(4)=f(2)$ : $2$ et $4$ sont symétriques par rapport à $3$, axe de symétrie de la parabole.

  \item $f(0)=-10$, donc $A(0\,;-10)$ appartient à la courbe.
  $f(1)=0$, donc $B(1\,;0)$ appartient à la courbe (c'est même une des deux racines).

  \item \textbf{Bilan.} Le signe de $k$ donne l'orientation de la parabole et le sens des limites ;
  $\alpha$ donne l'abscisse du sommet et l'axe de symétrie $x=\alpha$ ; $\beta$ donne l'ordonnée du sommet,
  c'est-à-dire la valeur de l'extremum (minimum si $k>0$, maximum si $k<0$). Le tableau de signes se déduit
  de la position de $\beta$ par rapport à $0$ et du signe de $k$.
\end{enumerate}

\begin{corrbox}
\textbf{Contrôle avec le document interactif.} Pour chaque partie, régler les curseurs sur les valeurs
de l'énoncé ($a=2$ ; $c=-3$ ; $a=-1$, $c=-4$ ; $k=2$ puis $k=-0{,}5$ ; $k=-2$, $\alpha=3$, $\beta=8$)
et vérifier : position du sommet, tableau de variations, tableau de signes, limites et calculs d'images.
\end{corrbox}

\end{document}